\section{Metodologia}

% TODO: redactar en pasado y en tercera persona, siguiendo el orden del
% pipeline implementado.

\subsection{Datos}

% TODO: COIL-100 \parencite{nene1996} y Fruits-360 \parencite{muresan2018}.
% Justificar el subconjunto de diez clases con solapamiento deliberado y el
% balanceo a trescientas imagenes por clase.

\subsection{Segmentacion y normalizacion}

% TODO: umbral sobre la distancia al color de fondo en YCbCr, operaciones
% morfologicas, contorno de area maxima, recorte cuadrado con fondo neutro.

\subsection{Representaciones evaluadas}

% TODO: via A (dieciseis descriptores geometricos y cromaticos), via B
% (eigen-objetos sobre imagenes de 64x64), via C (embeddings de
% MobileNetV3-Small congelado \parencite{howard2019} reducidos por PCA propia).

\subsection{Clasificador}

% TODO: minimos cuadrados multiclase con regularizacion de Tikhonov;
% estandarizacion ajustada unicamente con el conjunto de entrenamiento.

\subsection{Protocolo experimental}

% TODO: particion estratificada con semilla fija, validacion cruzada de cinco
% pliegues, tres barridos de ablacion y baselines de comparacion.

\subsection{Implementacion}

% TODO: NumPy como unico motor; sklearn empleado solo como referencia de
% verificacion numerica; torch empleado solo como extractor congelado.
